\section{Effective representation as a compatibility audit}
\label{sec:effective-representation-audit}

The role of effective-field language in the present construction is best
illustrated by the curvature-aware scalar perturbation audit performed during
the development of the R1 cosmological bridge.  A conventional
Newtonian-gauge effective representation was constructed in order to compare
the geometry-first response with familiar late-time perturbation variables.
After index and velocity-convention corrections, the system still exhibited a
curvature-dependent global constraint-preservation defect.  The corresponding
flattened system closed numerically, while the curved system did not satisfy
the same preservation criterion.

This result is not interpreted as a general failure of effective field theory.
Nor is it used as evidence against quantum field theory.  It establishes a
more limited point: an effective variable set and its equations can provide a
useful local or phenomenological representation without necessarily being a
globally faithful realization of the complete constrained geometry.

The response of this work was therefore not to modify the effective system
until it reproduced the desired cosmological signal.  Instead, the defining
transport was returned upstream to the physical action reduction:
\begin{equation}
\boxed{
\text{physical Hessian}
\rightarrow
X_2
\rightarrow
\mathcal R_{g2}X_2
\rightarrow
\text{optical transport}
\rightarrow
\text{observable kernel}.
}
\end{equation}
The curvature-aware effective equations remain in the project as a
compatibility audit.  They are not the ontological definition of the BHSM
mode.

\subsection{A concrete methodological distinction}

The practical difference may be summarized as follows:
\begin{center}
\begin{tabular}{p{0.43\linewidth} p{0.48\linewidth}}
\textbf{Representation-first question} &
\textbf{Geometry-first question}\\
\hline
Which fields should be introduced? &
Which perturbations survive the physical constraint and gauge quotient?\\
Which coefficients fit the observations? &
Which coefficients are matrix elements of the retained action?\\
Is the effective equation stable? &
Is the physical Hessian admissible on the common domain?\\
Can the field reproduce the signal? &
Does one reduced state predict multiple observables consistently?\\
Can missing behavior be absorbed by new terms? &
Does a failed prediction reject the proposed reduction?\\
\end{tabular}
\end{center}

This difference becomes scientifically meaningful only through prediction and
failure.  Several candidate reductions tested in this work fail that
criterion.  A homogeneous two-state mode cannot reproduce the required
low-redshift localization.  Preregistered foreground number-density,
luminosity-weighted, and seam-weighted line-of-sight proxies do not survive
held-out transfer.  These failures are retained as part of the evidence
because the rejected reductions were not subsequently tuned to recover the
signal.

\subsection{Fields as response coordinates}

The geometry-first interpretation does not prohibit field variables.  It
changes their logical status.  Once constraints, boundary data, gauge
redundancy, and nonselected directions are removed, a field-like variable may
serve as an efficient coordinate on the physical response.  Different
effective coordinate choices may therefore represent the same underlying
reduced geometry.

This observation motivates the optical formulation developed in this work.
A photon does not require an object-by-object effective-field inventory of
everything along its path.  It responds to the total spacetime geometry
through the optical tidal matrix.  The Jacobi/Sachs system consequently
provides a direct bridge from the reduced geometry to observables without
assigning fundamental status to every intermediate field representation.
